\section{Integral dependence}
\label{am-ch5-integral}
\begin{definition}[Integral element and integral extension]
\label{am-def-integral}
\source{atiyah-macdonald-commutative-algebra:page-0068}
For rings $A\subseteq B$, an element $x\in B$ is integral over $A$ if it satisfies
$x^n+a_1x^{n-1}+\cdots+a_n=0$ with $a_i\in A$.  The integral closure of $A$
in $B$ is the set of all such elements; $A$ is integrally closed in $B$ when
this equals $A$, and $B$ is integral over $A$ when every element is integral.
\end{definition}
\begin{proposition}[Equivalent tests for integrality]
\label{am-prop-5-1}
\dcref{5.1}
\source{atiyah-macdonald-commutative-algebra:page-0068}
For $x\in B$, the following are equivalent: (i) $x$ is integral over $A$;
(ii) $A[x]$ is finite as an $A$-module; (iii) $A[x]$ lies in a subring finite
over $A$; (iv) some faithful $A[x]$-module is finite over $A$.
\end{proposition}
\begin{proof}
An integral equation expresses every power of $x$ in the span of
$1,x,\ldots,x^{n-1}$, proving (i)$\Rightarrow$(ii)$\Rightarrow$(iii)$\Rightarrow$(iv)
by taking the subring itself.  For (iv), multiplication by $x$ on the finite
$A$-module satisfies the determinant trick (Proposition~\ref{am-prop-2-4});
faithfulness turns the resulting endomorphism equation into the required
monic polynomial for $x$.
\uses{am-prop-2-4}
\end{proof}
\begin{corollary}[Finite algebra generated by integral elements]
\label{am-cor-5-2}
\dcref{5.2}
\source{atiyah-macdonald-commutative-algebra:page-0069}
If $x_1,\ldots,x_n$ are integral over $A$, then $A[x_1,\ldots,x_n]$ is finite
over $A$.
\uses{am-prop-5-1,am-prop-2-16}
\end{corollary}
\begin{proof}
Induct on $n$.  Adjoin $x_n$ to the finite algebra generated by the preceding
elements and use Proposition~\ref{am-prop-5-1}, then transitivity of finite
generation (Proposition~\ref{am-prop-2-16}).
\end{proof}
\begin{corollary}[Integral closure is a ring]
\label{am-cor-5-3}
\dcref{5.3}
\source{atiyah-macdonald-commutative-algebra:page-0069}
The elements integral over $A$ form a subring containing $A$.
\uses{am-cor-5-2,am-prop-5-1}
\end{corollary}
\begin{proof}
For integral $x,y$, the finite $A$-module $A[x,y]$ contains $x+y$ and $xy$;
Proposition~\ref{am-prop-5-1}(iii) makes both integral.
\end{proof}
\begin{corollary}[Transitivity]
\label{am-cor-5-4}
\dcref{5.4}
\source{atiyah-macdonald-commutative-algebra:page-0069}
If $A\subseteq B\subseteq C$, $B$ is integral over $A$, and $C$ integral over
$B$, then $C$ is integral over $A$.
\end{corollary}
\begin{proof}
For $x\in C$, its integral equation has finitely many coefficients in $B$.
The $A$-algebra they generate is finite by Corollary~\ref{am-cor-5-2}; adjoining
$x$ is finite over it, hence finite over $A$, and Proposition~\ref{am-prop-5-1}
applies.
\uses{am-cor-5-2,am-prop-2-16,am-prop-5-1}
\end{proof}
\begin{corollary}[The integral closure is integrally closed]
\label{am-cor-5-5}
\dcref{5.5}
\source{atiyah-macdonald-commutative-algebra:page-0070}
The integral closure of $A$ in $B$ is integrally closed in $B$.
\uses{am-cor-5-4}
\end{corollary}
\begin{proof}
An element integral over the closure is integral over $A$ by transitivity, hence
belongs to the closure.
\end{proof}
\begin{proposition}[Quotients and fractions of integral extensions]
\label{am-prop-5-6}
\dcref{5.6}
\source{atiyah-macdonald-commutative-algebra:page-0070}
If $B$ is integral over $A$, then $B/\bideal$ is integral over
$A/(\bideal\cap A)$ for every ideal $\bideal\subseteq B$, and $S^{-1}B$ is
integral over $S^{-1}A$ for every multiplicative $S\subseteq A$.
\end{proposition}
\begin{proof}
Reduce an integral equation modulo $\bideal$.  For $x/s\in S^{-1}B$, divide
the equation for $x$ by the appropriate power of $s$ to obtain a monic equation
over $S^{-1}A$.
\end{proof}
\begin{proposition}[Integral domains and fields]
\label{am-prop-5-7}
\dcref{5.7}
\source{atiyah-macdonald-commutative-algebra:page-0070}
If integral domains $A\subseteq B$ satisfy that $B$ is integral over $A$, then
$B$ is a field iff $A$ is a field.
\end{proposition}
\begin{proof}
If $A$ is a field, use a minimal integral equation for $0\ne y\in B$ and
solve for $y^{-1}$.  If $B$ is a field and $0\ne x\in A$, apply an integral
equation to $x^{-1}\in B$ and multiply by a suitable power of $x$ to express
$x^{-1}$ in $A$.
\end{proof}
\begin{corollary}[Maximal ideals in integral extensions]
\label{am-cor-5-8}
\dcref{5.8}
\source{atiyah-macdonald-commutative-algebra:page-0070}
If $B$ is integral over $A$ and $\q\in\Spec B$ contracts to $\p$, then $\q$ is
maximal iff $\p$ is maximal.
\uses{am-prop-5-6,am-prop-5-7}
\end{corollary}
\begin{proof}
Apply Proposition~\ref{am-prop-5-7} to the integral domain extension
$A/\p\subseteq B/\q$.
\end{proof}
\begin{corollary}[Incomparability]
\label{am-cor-5-9}
\dcref{5.9}
\source{atiyah-macdonald-commutative-algebra:page-0070}
If $\q\subseteq\q'$ are primes of an integral extension $B/A$ with equal
contractions, then $\q=\q'$.
\uses{am-prop-5-6,am-cor-5-8,am-cor-3-13}
\end{corollary}
\begin{proof}
Pass modulo the common contraction and localize there.  The two extended
ideals in the resulting integral extension of a local ring are both maximal by
Corollary~\ref{am-cor-5-8}, so equal; contraction back gives the claim.
\end{proof}
\begin{theorem}[Lying over]
\label{am-thm-5-10}
\dcref{5.10}
\source{atiyah-macdonald-commutative-algebra:page-0071}
If $B$ is integral over $A$, every prime $\p$ of $A$ is the contraction of a
prime of $B$.
\end{theorem}
\begin{proof}
By Proposition~\ref{am-prop-5-6}, $B_\p$ is integral over $A_\p$.  Choose a
maximal ideal of $B_\p$; its contraction is the maximal ideal of the local
ring $A_\p$.  Contracting back to $B$ gives a prime over $\p$.
\uses{am-prop-5-6,am-cor-5-8,am-thm-1-3}
\end{proof}
\begin{theorem}[Going up]
\label{am-thm-5-11}
\dcref{5.11}
\source{atiyah-macdonald-commutative-algebra:page-0071}
For an integral extension $A\subseteq B$, a chain of primes in $A$ extending a
given chain in $B$ can be lifted one step at a time: if
$\p_1\subseteq\cdots\subseteq\p_n$ and
$\q_1\subseteq\cdots\subseteq\q_m$ contract correspondingly, $m<n$, there
are $\q_{m+1},\ldots,\q_n$ with $\q_i\cap A=\p_i$.
\end{theorem}
\begin{proof}
Induct to the one-step case.  Modulo $\q_1$ and $\p_1$, the extension remains
integral; lying over gives a prime over the image of $\p_2$, whose inverse image
is the required next prime.  Repeat.
\uses{am-thm-5-10,am-prop-5-6}
\end{proof}

\section{Integral closure and going down}
\label{am-ch5-going-down}
\begin{proposition}[Localization of integral closure]
\label{am-prop-5-12}
\dcref{5.12}
\source{atiyah-macdonald-commutative-algebra:page-0071}
If $C$ is the integral closure of $A$ in $B$, then $S^{-1}C$ is the integral
closure of $S^{-1}A$ in $S^{-1}B$.
\end{proposition}
\begin{proof}
One inclusion is Proposition~\ref{am-prop-5-6}.  If $b/s$ is integral after
localization, clear denominators in its monic equation; a suitable multiple
$bt$ is integral over $A$, so $b/s=(bt)/(st)\in S^{-1}C$.
\uses{am-prop-5-6}
\end{proof}
\begin{definition}[Integrally closed domain]
\label{am-def-integrally-closed}
\source{atiyah-macdonald-commutative-algebra:page-0072}
An integral domain is integrally closed when it is integrally closed in its field
of fractions.
\end{definition}
\begin{proposition}[Integral closedness is local]
\label{am-prop-5-13}
\dcref{5.13}
\source{atiyah-macdonald-commutative-algebra:page-0072}
For a domain $A$, integrally closedness is equivalent to integrally closedness of
all $A_\p$, and equivalently all $A_\m$.
\end{proposition}
\begin{proof}
Let $C$ be the integral closure of $A$ in its fraction field.  By
Proposition~\ref{am-prop-5-12}, $A_\p$ is integrally closed exactly when the
localized map $A\to C$ is surjective.  Proposition~\ref{am-prop-3-9} detects
surjectivity at all primes or maximals.
\uses{am-prop-5-12,am-prop-3-9}
\end{proof}
\begin{definition}[Integral over an ideal]
\label{am-def-integral-ideal}
\source{atiyah-macdonald-commutative-algebra:page-0072}
An element is integral over an ideal $\aideal\subseteq A$ if it satisfies a
monic equation whose coefficients all lie in $\aideal$.
\end{definition}
\begin{lemma}[Closure over an ideal]
\label{am-lem-5-14}
\dcref{5.14}
\source{atiyah-macdonald-commutative-algebra:page-0072}
If $C$ is the integral closure of $A$ in $B$ and $\aideal^e$ is the extension in
$C$, the elements of $B$ integral over $\aideal$ are exactly
$\rad(\aideal^e)$.
\end{lemma}
\begin{proof}
An integral equation over $\aideal$ gives $x^n\in\aideal^e$.  Conversely,
$x^n=\sum a_ix_i$ with $a_i\in\aideal,x_i\in C$; the finite module generated
by the $x_i$ is stable under multiplication by $x^n$ into $\aideal$, so the
determinant trick gives an integral equation over $\aideal$.
\uses{am-prop-2-4,am-cor-5-2,am-def-integral-ideal}
\end{proof}
\begin{proposition}[Minimal polynomial coefficients]
\label{am-prop-5-15}
\dcref{5.15}
\source{atiyah-macdonald-commutative-algebra:page-0072}
If $A$ is integrally closed, $A\subseteq B$ are domains, and $x\in B$ is
integral over an ideal $\aideal$, then the coefficients of the minimal polynomial
of $x$ over $\Frac(A)$ lie in $\rad(\aideal)$.
\end{proposition}
\begin{proof}
All conjugates of $x$ satisfy the same integral equation over $\aideal$; their
elementary symmetric functions, the minimal-polynomial coefficients, are
integral over $\aideal$ by Lemma~\ref{am-lem-5-14}.  Since $A$ is integrally
closed, they lie in $\rad(\aideal)$.
\uses{am-lem-5-14,am-def-integrally-closed}
\end{proof}
\begin{theorem}[Going down]
\label{am-thm-5-16}
\dcref{5.16}
\source{atiyah-macdonald-commutative-algebra:page-0073}
If $A$ is an integrally closed domain and $B$ is integral over $A$, every finite
descending chain of primes in $A$ extending a chain in $B$ can be lifted one
step at a time.
\end{theorem}
\begin{proof}
Reduce to $\p_1\supseteq\p_2$ and localize at $\q_1$.  It suffices by
Proposition~\ref{am-prop-3-16} to show the contraction of $B_{\q_1}\p_2$ is
$\p_2$.  If $x$ lies in this contraction, write $x=y/s$ with $y$ integral over
$\p_2$ and $s\notin\q_1$.  Proposition~\ref{am-prop-5-15} gives a minimal
equation whose coefficients lie in $\p_2$; applying it to $s=yx^{-1}$ and
using integrality of $s$ forces $x\in\p_2$, a contradiction otherwise.  Thus
the contraction is $\p_2$, and Proposition~\ref{am-prop-3-16} supplies the
prime below $\q_1$.
\uses{am-prop-3-16,am-prop-5-15,am-prop-3-11}
\end{proof}
\begin{proposition}[Finite integral closure in a separable extension]
\label{am-prop-5-17}
\dcref{5.17}
\source{atiyah-macdonald-commutative-algebra:page-0073}
If $A$ is integrally closed with fraction field $K$, $L/K$ is finite separable,
and $B$ is the integral closure of $A$ in $L$, then a $K$-basis
$v_1,\ldots,v_n$ of $L$ exists with $B\subseteq\sum_iAv_i$.
\end{proposition}
\begin{proof}
Scale any basis so its elements lie in $B$.  Use the nondegenerate trace pairing
to choose a dual basis $v_i$ with $\operatorname{Tr}(u_iv_j)=\delta_{ij}$.  For
$x=\sum x_iv_i\in B$, each $\operatorname{Tr}(xu_j)$ lies in $A$ by the
integral equation and integrally closedness, and equals $x_j$; hence $x_i\in A$.
\uses{am-def-integrally-closed}
\end{proof}

\section{Valuation rings}
\label{am-ch5-valuations}
\begin{definition}[Valuation ring]
\label{am-def-valuation-ring}
\source{atiyah-macdonald-commutative-algebra:page-0074}
For a domain $B$ with fraction field $K$, $B$ is a valuation ring of $K$ if for
every nonzero $x\in K$, either $x\in B$ or $x^{-1}\in B$.
\end{definition}
\begin{proposition}[Basic valuation properties]
\label{am-prop-5-18}
\dcref{5.18}
\source{atiyah-macdonald-commutative-algebra:page-0074}
A valuation ring is local; every intermediate ring $B\subseteq B'\subseteq K$
is a valuation ring; and a valuation ring is integrally closed.
\end{proposition}
\begin{proof}
The nonunits form an ideal: multiplication by an element outside it cannot
turn a nonunit into a unit, and for nonzero nonunits $x,y$ one of $x/y,y/x$
lies in $B$, giving $x+y\in\mathfrak m$.  The valuation condition is inherited
by intermediate rings.  If $x$ is integral and $x^{-1}\in B$, divide its monic
equation by the appropriate power of $x$ to express $x\in B$.
\end{proof}
\begin{lemma}[Maximal valued subring is local]
\label{am-lem-5-19}
\dcref{5.19}
\source{atiyah-macdonald-commutative-algebra:page-0074}
For a maximal pair $(B,g)$ consisting of a subring $B\subseteq K$ and a map to
an algebraically closed field, $B$ is local and $\Ker g$ is its maximal ideal.
\end{lemma}
\begin{proof}
The kernel is prime.  The map extends to $B_{\Ker g}$ because denominators map
to nonzero elements.  Maximality forces $B=B_{\Ker g}$, which is local.
\end{proof}
\begin{lemma}[One of $x,x^{-1}$ can be adjoined]
\label{am-lem-5-20}
\dcref{5.20}
\source{atiyah-macdonald-commutative-algebra:page-0074}
For $0\ne x\in K$, if the extensions of the maximal ideal to both $B[x]$ and
$B[x^{-1}]$ were unit ideals, then minimal-degree relations in $x$ and $x^{-1}$
contradict one another.  Hence one extension is proper.
\end{lemma}
\begin{proof}
Write $1=\sum u_ix^i$ and $1=\sum v_jx^{-j}$ with coefficients in the maximal
ideal, choosing the largest exponents minimally.  Multiply the second relation
by the larger power of $x$ and use that $1-v_0$ is a unit (Lemma~\ref{am-lem-5-19})
to eliminate the highest term in the first, contradicting minimality.
\uses{am-lem-5-19}
\end{proof}
\begin{theorem}[Existence of valuation rings]
\label{am-thm-5-21}
\dcref{5.21}
\source{atiyah-macdonald-commutative-algebra:page-0075}
Every maximal pair $(B,g)$ as above has $B$ a valuation ring of $K$.
\end{theorem}
\begin{proof}
For $x\ne0$, Lemma~\ref{am-lem-5-20} lets us assume the maximal ideal of
$B[x]$ is proper.  Put it inside a maximal ideal and reduce; the residue field
is a finite algebraic extension of the residue field of $B$, so the embedding
extends into the algebraically closed target.  Maximality then forces $B[x]=B$,
so $x\in B$ (or apply the argument to $x^{-1}$).
\uses{am-lem-5-19,am-lem-5-20}
\end{proof}
\begin{corollary}[Integral closure as intersection of valuation rings]
\label{am-cor-5-22}
\dcref{5.22}
\source{atiyah-macdonald-commutative-algebra:page-0075}
The integral closure of a subring $A$ in a field $K$ is the intersection of all
valuation rings of $K$ containing $A$.
\uses{am-prop-5-18,am-thm-5-21,am-thm-1-3}
\end{corollary}
\begin{proof}
Every such valuation ring is integrally closed.  If $x$ is outside the integral
closure, apply Theorem~\ref{am-thm-5-21} to a maximal ideal of $A[x^{-1}]$
containing $x^{-1}$; the resulting valuation ring contains $A$ but not $x$.
\end{proof}
\begin{proposition}[Specialization avoiding a nonzero element]
\label{am-prop-5-23}
\dcref{5.23}
\source{atiyah-macdonald-commutative-algebra:page-0075}
If $A\subseteq B$ are domains, $B$ is finite type over $A$, and $0\ne v\in B$,
there is $0\ne u\in A$ such that every map $A\to\Omega$ to an algebraically
closed field with $u$ nonzero extends to $B\to\Omega$ sending $v$ nonzero.
\end{proposition}
\begin{proof}
Induct on algebra generators.  For one transcendental generator, choose the
leading coefficient of a polynomial expression for $v$ and avoid its finitely
many roots in the infinite field.  For an algebraic generator, clear the
leading coefficients of integral equations for the generator and $v^{-1}$;
after inverting their product, Theorem~\ref{am-cor-5-22} places both in a
valuation ring, where $v$ is a unit.  Restrict its map to $B$.
\uses{am-cor-5-22}
\end{proof}
\begin{corollary}[Finite algebra field extensions]
\label{am-cor-5-24}
\dcref{5.24}
\source{atiyah-macdonald-commutative-algebra:page-0076}
If $k$ is a field and a finitely generated $k$-algebra $B$ is a field, then $B$
is a finite algebraic extension of $k$.
\uses{am-prop-5-23}
\end{corollary}
\begin{proof}
Apply Proposition~\ref{am-prop-5-23} with $A=k$ and $v=1$.  If $B$ had a
transcendental generator, a map to an algebraic closure would extend, producing
a nontrivial $k$-homomorphism from the field $B$ to the algebraic closure and
forcing all generators algebraic.  A finite set of algebraic generators gives a
finite-dimensional extension by Corollary~\ref{am-cor-5-2}.
\end{proof}
